\documentclass[10pt,a4paper]{amsart}
\usepackage[margin=19mm]{geometry}
\usepackage{amsmath,amssymb,mathtools}
\usepackage[scr=rsfs,cal=euler]{mathalfa}
\usepackage{xfrac}
\usepackage[unicode,hidelinks,bookmarks=false]{hyperref}
\newcommand{\ie}{i.e.\ }
\newcommand{\cf}{cf.\ }
\newcommand{\resp}{respectively\ }
\newcommand{\Subsection}{\subsection}
\providecommand{\nobreakdash}{\nobreak\hbox{-}}
\setlength{\emergencystretch}{2em}
\multlinegap0pt
\usepackage{tikz}
\newtheorem{prop}{Proposition}
\DeclareMathOperator{\Aut}{Aut}
\DeclareMathOperator{\id}{id}
\DeclareMathOperator{\ord}{ord}
\title{An introduction to quantum symmetries}
\author{Christian Voigt}
\date{Source: arXiv:2603.25141v1 (CC BY 4.0)}
\begin{document}
\maketitle

\noindent\emph{Source and licence.} An introduction to quantum symmetries. Version arXiv:2603.25141v1 (CC BY 4.0), distributed by arXiv under the Creative Commons Attribution 4.0 International licence, http://creativecommons.org/licenses/by/4.0/, as recorded in the arXiv licence metadata for this exact version and shown on the abstract page https://arxiv.org/abs/2603.25141v1. Full licence notice: \texttt{licenses/arxiv-2603.25141-CC-BY-4.0.txt}.\par\smallskip

\noindent\textit{Editorial scope.} Counted unit: the article's prop disjointautomorphisms (source0.tex lines 1085--1089). The article's complete original proof of it is delivered with it (source0.tex lines 1091--1141, one \texttt{proof} environment of 51 lines). No mathematics is written, repaired or re-derived: the statement and the proof are the article's own lines, in the article's own notation, and no retained line is substituted or re-flowed.  No further source-local context is needed: every cross-reference of every retained line resolves inside the two delivered spans.  Nothing was omitted from any retained span, so the package carries no omission record.  No retained line contains a citation, so the wrapper carries no bibliography and no bibliography entry.  No retained line contains a figure environment, an image-inclusion command or drawing code, so no image is delivered and the package declares no asset dependency.  Editorial formatting only, in addition: an AMS \texttt{amsart} preamble that restates the article's own declarations and macros used by the retained lines, namely the environments \texttt{prop} on separate counters, together with the standard packages those lines need.  The retained environments print in this document as Proposition 1 in order of appearance, whereas the article prints the same items with the numbering of its own full document.  The complete native source of the article as distributed in the arXiv e-print for this exact version is bundled unchanged as \texttt{sources/197189/source0.tex}.
\begin{prop} \label{disjointautomorphisms}
Let $ X $ be a graph admitting a pair of disjoint automorphisms $ \sigma, \tau \in \Aut(X) $, and assume that $ k \in \mathbb{N} $ does not exceed the order
of neither of these automorphisms. Then $ X $ admits an irreducible quantum automorphism of dimension $ k $. In particular, if $ \Aut(X) $ contains a pair of nontrivial 
disjoint automorphisms then $ X $ has quantum symmetry. 
\end{prop} 

\begin{proof} 
Let us fix $ k \in \mathbb{N} $ such that $ k \leq \min(\ord(\sigma), \ord(\tau)) $, and note that if $ \ord(\sigma) = \infty = \ord(\tau) $ this means that we can 
choose $ k $ arbitrarily. We shall construct an irreducible quantum automorphism $ \rho = (\mathbb{C}^k, u^\rho) $ as follows. 

Let $ G = \mathbb{Z}/k\mathbb{Z} $ and consider the actions of $ C(G) $ and $ C^*(G) $ on $ \mathbb{C}^k = l^2(G) $, induced by pointwise multiplication of functions in $ C(G) $, 
and the regular representation of $ G $, respectively. We shall write $ p_1, \dots p_k $ and $ q_1 \dots, q_k $ for the images in $ M_k(\mathbb{C}) = B(l^2(G)) $ of the minimal 
projections in $ C(G) $ and $ C^*(G) $ under these representations, respectively. 
 
Define $ u^\rho = (u^\rho_{x,y})_{x,y \in V_X} $ by $ u^\rho = \sum_{r = 1}^k u^{\sigma^r} p_r + \sum_{s = 1}^k u^{\tau^s} q_s - u^{\id} 1 $. That 
is, the matrix $ u^\rho_{x,y} \in M_k(\mathbb{C}) $ is given by 
$$ 
u^\rho_{x,y} = \sum_{r = 1}^k \delta_{x,\sigma^r(y)} p_r + \sum_{s = 1}^k \delta_{x,\tau^s(y)} q_s - \delta_{x,y} 1
$$ 
for $ x, y \in V_X $. 
Since $ \sigma $ and $ \tau $ are graph automorphisms we clearly have $ u^\rho A_X = A_X u^\rho $. Let 
\begin{align*}
M_{xy} &= \{1 \leq r \leq k \mid \sigma^r(y) = x \}, \\
N_{xy} &= \{1 \leq s \leq k \mid \tau^s(y) = x \}, 
\end{align*}
and observe that 
$$ 
u^\rho_{x,y} = 
\begin{cases} 
\sum_{r \in M_{xy}} p_r & \text{ if } \sigma(y) \neq y  \\
\sum_{s \in N_{xy}} q_s & \text{ if } \tau(y) \neq y  \\
\delta_{x,y} 1 & \text{ if } \sigma(y) = y = \tau(y). 
\end{cases}
$$
In particular, every $ u^\rho_{x,y} $ for $ x, y \in V_X $ is a projection. In addition, we have 
\begin{align*}
\sum_{x \in V_X} u^\rho_{x,y} &= \sum_{r = 1}^k \sum_{x \in V_X} \delta_{x, \sigma^r(y)} p_r + \sum_{s = 1}^k \sum_{x \in V_X} \delta_{x, \tau^s(y)} q_s - 1
= \sum_{r = 1}^k  p_r + \sum_{s = 1}^k q_s - 1 = 1,  
\end{align*}
and similarly 
\begin{align*}
\sum_{y \in V_X} u^\rho_{x,y} &= \sum_{r = 1}^k \sum_{y \in V_X} \delta_{x,\sigma^r(y)} p_r + \sum_{s = 1}^k \sum_{y \in V_X} \delta_{x,\tau^s(y)} q_s - 1 
= \sum_{r = 1}^k p_r + \sum_{s = 1}^k q_s - 1 = 1 
\end{align*}
as required. It follows that $ \rho = (\mathbb{C}^k, u^\rho) $ is a quantum automorphism of $ X $. 

Let us now check that $ \rho $ is irreducible. If $ \ord(\sigma) = m < \infty $ then 
upon decomposing $ X $ into the orbits of $ \sigma $ we see that there are finitely many elements $ v_1, \dots, v_a \in V_X $, fixed under $ \tau $, 
such that $ (\sigma^t(v_1), \dots, \sigma^t(v_a)) = (v_1, \dots, v_a) $ for $ 0 \leq t < k $ implies $ t = 0 $. 
If $ \ord(\sigma) = \infty $ we either find an infinite orbit, or orbits of arbitrarily large finite size. 
Again this allows us to choose $ v_1, \dots, v_a \in V_X $ such that $ (\sigma^t(v_1), \dots, \sigma^t(v_a)) = (v_1, \dots, v_a) $ for $ 0 \leq t < k $ implies $ t = 0 $.
In fact, we may take $ v_1 = \cdots = v_a $ for a suitably chosen vertex in this case. 
In either case, it follows that $ u^\rho_{\sigma^r(v_1), v_1} \cdots u^\rho_{\sigma^r(v_a), v_a} = p_r $. 

In the same way we find $ w_1, \dots, w_b \in V_X $ such that $ u^\rho_{\tau^r(w_1), w_1} \cdots u^\rho_{\tau^r(w_b), w_b} = q_r $. 
Since the projections $ p_i, q_j $ for $ 1 \leq i,j \leq k $ generate $ M_k(\mathbb{C}) $ this yields the claim.
\end{proof} 

\end{document}
